\section{Inizializzazione e Firme degli Strumenti}
Il Listato \ref{lst:fastmcp-server} fornisce un estratto architetturalmente rilevante dell'implementazione del server Model Context Protocol (MCP) sviluppato per il framework Service-Forge. Il codice evidenzia esclusivamente l'inizializzazione del servizio, le istruzioni di base fornite al modello linguistico e la registrazione degli strumenti operativi (\emph{tools}) tramite tipizzazione forte, la quale agisce da schema di validazione statica per i payload JSON. Per questioni di leggibilità, la logica procedurale di basso livello è stata omessa.

\begin{lstlisting}[language=Python, label={lst:fastmcp-server}]
import os
import sys
from typing import Any, Dict, List, Optional

from mcp.server.fastmcp import FastMCP
from cli.profile_setup import ProfileSetup
from core.engine import GenerationEngine

mcp = FastMCP(
    "service-forge",
    instructions=(
        "Service Forge generates AWS Lambda-based microservices. "
        "Use get_profiles to discover available profiles and services. "
        "IMPORTANT: Every AWS service has configurable parameters (bucket names, "
        "table names, secret names) that you MUST collect from the user. "
        "Never use default placeholder values without asking the user first. "
        "Use create_session to start building a configuration, then configure "
        "properties with configure_service."
    ),
)

# ==========================================================
# Resources - Dati in sola lettura ispezionabili dall'LLM
# ==========================================================

@mcp.resource("forge://profiles")
def resource_profiles() -> str:
    """All available profiles with metadata."""
    # [... LOGICA PROCEDURALE OMESSA ...]
    pass

@mcp.resource("forge://aws-services")
def resource_aws_services() -> str:
    """All AWS services with their configurable keys, defaults, and descriptions."""
    # [... LOGICA PROCEDURALE OMESSA ...]
    pass

# ==========================================================
# Tools - Azioni invocabili dall'agente LLM
# ==========================================================

@mcp.tool()
def check_prerequisites() -> str:
    """
    Check if all prerequisites for project generation are met.
    Call this before generate_project to ensure the environment is ready.
    """
    # [... LOGICA PROCEDURALE OMESSA ...]
    pass

@mcp.tool()
def get_profiles() -> str:
    """
    List all active Service Forge profiles and available AWS services.
    Call this first to understand what can be generated.
    """
    # [... LOGICA PROCEDURALE OMESSA ...]
    pass

@mcp.tool()
def create_session() -> str:
    """
    Start a new project configuration session.
    Returns a session_id to use in subsequent tool calls.
    """
    # [... LOGICA PROCEDURALE OMESSA ...]
    pass

@mcp.tool()
def configure_service(
    session_id: str,
    service_name: str,
    config_values: Optional[Dict[str, str]] = None,
) -> str:
    """
    Add or configure an AWS service for the project.
    IMPORTANT: Always ask the user for the actual values.
    """
    # [... LOGICA PROCEDURALE OMESSA ...]
    pass

@mcp.tool()
def validate_config(session_id: str) -> str:
    """
    Validate the current session configuration.
    Returns errors if any, or confirms the config is ready for generation.
    """
    # [... LOGICA PROCEDURALE OMESSA ...]
    pass

@mcp.tool()
def generate_project(session_id: str) -> str:
    """
    Generate the project from the current session configuration.
    Creates files on disk and initializes a git repository.
    """
    # [... LOGICA PROCEDURALE OMESSA ...]
    pass

@mcp.tool()
def quick_generate(
    profile: str,
    app_name: str,
    project_path: str,
    aws_services: Optional[Dict[str, Dict[str, str]]] = None,
    included_services: Optional[List[str]] = None,
    runtime: Optional[Dict[str, Any]] = None,
) -> str:
    """
    One-shot project generation without session management.
    Builds config, validates, and generates in a single call.
    """
    # [... LOGICA PROCEDURALE OMESSA ...]
    pass

if __name__ == "__main__":
    mcp.run(transport="stdio")
\end{lstlisting}

\section{Esempio template Jinja}
Il Listato \ref{lst:template-compose} mostra un esempio dei template generativi basati sul motore Jinja2, in particolare si tratta del template usato per generare il file \texttt{docker-compose.yaml} in tutti i microservizi. La sintassi illustra la separazione netta tra la logica di iniezione dinamica delle variabili e la struttura statica del codice sorgente, garantendo la standardizzazione formale dei microservizi prodotti dal framework aziendale.

\begin{lstlisting}[label={lst:template-compose}]
services:
  {{ service_name }}:
    build:
      context: .
      target: dev
    platform: linux/arm64
    restart: on-failure
    environment:
      {% if app_name -%}
      APP_NAME: {{ app_name }}
      {%- else -%}
      APP_NAME: {{ service_name }}
      {%- endif %}
      {%- for key, value in environment.items() %}
      {{ key }}: {{ value }}
      {%- endfor %}
      {%- block aws_services_block %}{%- endblock %}
      {%- if use_localstack %}
      AWS_ENDPOINT_URL: http://localstack:4566
      AWS_ACCESS_KEY_ID: test
      AWS_SECRET_ACCESS_KEY: test
      AWS_DEFAULT_REGION: {{ aws_region }}
      AWS_LAMBDA_FUNCTION_TIMEOUT: {{ lambda_timeout }}
      AWS_LAMBDA_FUNCTION_MEMORY_SIZE: {{ lambda_memory_size }}
      {%- endif %}
    volumes:
      - ./app:/var/task
      - ~/.aws:/root/.aws:ro
    ports:
      - {{ service_port }}:8080
      - 5678:5678
    {% if profile == 'service' -%}
    command:
      sh -c "echo 'Please run the debugger to proceed (F5)' && python -Xfrozen_modules=off -m debugpy --listen 0.0.0.0:5678 --wait-for-client -m uvicorn main:app --host 0.0.0.0 --port 8080 --reload --log-level debug"
    {%- endif  %}

{% if use_localstack %}
  localstack:
      image: localstack/localstack
      ports:
          - "127.0.0.1:4566:4566"
          - "127.0.0.1:4510-4559:4510-4559"
      env_file:
          - .env.template
          - .env
      environment:
          DEBUG: ${DEBUG:-0}
      volumes:
          - "${LOCALSTACK_VOLUME_DIR:-./localstack_volume}:/var/lib/localstack"
          - "/var/run/docker.sock:/var/run/docker.sock:ro"
          - "./init-aws.sh:/etc/localstack/init/ready.d/init-aws.sh:ro"
{%- endif%}
\end{lstlisting}